% ---
% id: meat-core/thm:oracle
% kind: thm
% title: "Limit of the value as the Rényi order tends to one"
% summary: >-
%   If every acceptable distribution gives each key-live value probability at
%   least some p₀ > 0 and some feasible state reproduces an acceptable
%   distribution, the value tends, as α → 1, to a blockwise relative-entropy
%   program. The limit is a validation reference and not a bound.
% prominence: standard
% status: implemented
% uses: [thm:branch-formula, lem:b-derivative, prop:exact-cancellation, thm:reduction, def:admissible-declaration, def:admissible-acceptance-set, def:feasible, foundations/lem:kernel, foundations/def:kl]
% mentions: [rem:oracle]
% assumes: []
% refs: []
% provenance: newly-derived
% implements: certiqs_key.meat_core.objective.limit:check_limit_hypotheses
% lean:
% ---

As the R\'enyi order tends to one, the weight $\alpha/(\alpha-1)$ of the
divergence term of the objective diverges, and the entropy part of each
branch tends to a relative entropy to its pinching. Remark~\ref{rem:oracle} says what the
limit validates and on which declarations it is available.

\begin{theorem}\label{thm:oracle}
Let $\Decl$ satisfy Definition~\ref{def:admissible-declaration} and let
$\acSet$ satisfy Definition~\ref{def:admissible-acceptance-set}. Assume that
there is $p_0>0$ with $q(\hat c)\ge p_0$ for every $\mathbf q\in\acSet$ and
$\hat c\in\Chkey$, and that the asymptotic feasible set
$\Feas_1\coloneqq\{\omega\in\Feas:\Phimap[\omega]\in\acSet\}$ is
nonempty. Then
\begin{equation}\label{eq:oracle}
  \lim_{\alpha\downarrow1}\hlow(\alpha)
  \;=\;h_1\coloneqq
  \inf_{\omega\in\Feas_1}\ \sum_{\hat c\in\Chkey}\
  \sum_{s\in\SecSet_K(\hat c)}
  \Dtwo\Big(\cG_{\hat c,s}(\omega_s)\,\Big\|\,
  \cZ\big(\cG_{\hat c,s}(\omega_s)\big)\Big).
\end{equation}
\end{theorem}

\begin{proof}
\emph{Idea.} The upper bound evaluates the objective at a state that
reproduces an acceptable distribution exactly, where only the entropy part
survives, and passes to the limit with the derivative of
Lemma~\ref{lem:b-derivative} at $\alpha=1$. The lower bound takes
near-optimal points along $\alpha_n\downarrow1$, uses the divergence part
to force a limit point that reproduces its distribution, and bounds the
entropy part from below by the continuous majorant of that lemma.

\emph{Notation.} Write $\sigma_{\hat c,s}(\omega)\coloneqq\cG_{\hat c,s}(\omega_s)$ and,
for $P_{\hat c}(\omega)>0$,
$e_{\hat c}(\alpha,\omega)\coloneqq\frac{1}{1-\alpha}\log_2
\frac{K_{\hat c}(\omega)}{P_{\hat c}(\omega)}\ge0$.

\emph{Upper bound.} Fix $\omega\in\Feas_1$ and take
$\mathbf q\coloneqq\Phimap[\omega]\in\acSet$. All divergence terms of
$O_{\alpha}$ vanish, and $q(\hat c)=P_{\hat c}(\omega)\ge p_0$ for
$\hat c\in\Chkey$. Then
\begin{equation*}
  O_{\alpha}(\mathbf q,\omega)
  =\sum_{\hat c\in\Chkey}P_{\hat c}\,e_{\hat c}(\alpha,\omega)
  =\sum_{\hat c\in\Chkey}\frac{P_{\hat c}}{(1-\alpha)\ln2}\,
  \ln\frac{K_{\hat c}(\alpha)}{P_{\hat c}} ,
\end{equation*}
and $K_{\hat c}(\alpha)-P_{\hat c}
=\sum_{s}\big(B_{\sigma_{\hat c,s}}(\alpha)
-B_{\sigma_{\hat c,s}}(1)\big)$ by \eqref{eq:c5-sector} and
$B_{\sigma}(1)=\tr\sigma$ together with
$P_{\hat c}-\sum_s\tr\sigma_{\hat c,s}=\tr\Theta_{\hat c}$ cancelling
between $K$ and $P$ (both formulas from \eqref{eq:fK}). By
Lemma~\ref{lem:b-derivative} and l'H\^opital (or the $C^1$ mean value
theorem), as $\alpha\downarrow1$,
\begin{equation*}
  O_{\alpha}(\mathbf q,\omega)\longrightarrow
  -\frac{1}{\ln2}\sum_{\hat c\in\Chkey}\sum_{s}
  B'_{\sigma_{\hat c,s}}(1)
  =\sum_{\hat c\in\Chkey}\sum_{s}
  \Dtwo\big(\sigma_{\hat c,s}\|\cZ(\sigma_{\hat c,s})\big),
\end{equation*}
where the sum may be restricted to $s\in\SecSet_K(\hat c)$ because
key-inert sectors have $\sigma$ already $\cZ$-block-diagonal
(Proposition~\ref{prop:exact-cancellation}) and hence zero relative entropy. Therefore
$\limsup_{\alpha\downarrow1}\hlow(\alpha)\le
h_1$-objective of $\omega$, for every $\omega\in\Feas_1$, so
$\limsup\hlow(\alpha)\le h_1$.

\emph{Lower bound.} Let $\alpha_n\downarrow1$ and pick feasible
$(\mathbf q_n,\omega_n)$ with
$O_{\alpha_n}(\mathbf q_n,\omega_n)\le\hlow(\alpha_n)+\frac1n$; by the
upper bound this sequence of values is bounded above by some $C<\infty$.
Decompose $O_{\alpha}=\frac{\alpha}{\alpha-1}\sfD(\mathbf q\|\Phimap)
+\sum_{\hat c\in\Chkey}q(\hat c)e_{\hat c}$ (undo
\eqref{eq:rearrange}); both parts are non-negative (the divergence by
Gibbs' inequality of \xref{foundations/def:kl}, $\mathbf q_n$ and
$\Phimap[\omega_n]$ being probability vectors, and $e_{\hat c}\ge0$ since $K_{\hat c}\le P_{\hat c}$ by
Theorem~\ref{thm:branch-formula}'s \emph{cone clause}, the points $\omega_n$ being
feasible but not pure). Hence
$\sfD(\mathbf q_n\|\Phimap[\omega_n])\le C\frac{\alpha_n-1}{\alpha_n}
\to0$. By compactness pass to a subsequence with
$(\mathbf q_n,\omega_n)\to(\mathbf q^{*},\omega^{*})
\in\acSet\times\Feas$: $\acSet$ is a closed subset of the simplex
(Definition~\ref{def:admissible-acceptance-set}) and $\Feas$ is compact
(Definition~\ref{def:feasible}). $\sfD(\mathbf q\|\mathbf p)$ is the sum
over $\hat c$ of the kernel $d(q(\hat c),p(\hat c))$ of
\xref{foundations/lem:kernel}, which is lower semicontinuous, and
$\Phimap$ is continuous, so
$\sfD(\mathbf q^{*}\|\Phimap[\omega^{*}])\le
\liminf_n\sfD(\mathbf q_n\|\Phimap[\omega_n])=0$; by the equality case of
Gibbs' inequality (\xref{foundations/def:kl})
$\mathbf q^{*}=\Phimap[\omega^{*}]$ and
$\omega^{*}\in\Feas_1$; in particular
$P_{\hat c}(\omega^{*})=q^{*}(\hat c)\ge p_0$ for $\hat c\in\Chkey$, so
$P_{\hat c}(\omega_n)\ge p_0/2$ eventually. For each
$\hat c\in\Chkey$, using $\ln(1+u)\le u$ with
$u_n\coloneqq\frac{K_{\hat c}(\omega_n)-P_{\hat c}(\omega_n)}
{P_{\hat c}(\omega_n)}\in[-1,0]$ and $\frac{1}{1-\alpha_n}<0$:
\begin{equation*}
  e_{\hat c}(\alpha_n,\omega_n)
  \;\ge\;\frac{u_n}{(1-\alpha_n)\ln2}
  \;=\;\frac{1}{\ln2\,P_{\hat c}(\omega_n)}
  \sum_{s}\Big[-\frac{B_{\sigma_{\hat c,s}(\omega_n)}(\alpha_n)
  -B_{\sigma_{\hat c,s}(\omega_n)}(1)}{\alpha_n-1}\Big].
\end{equation*}
By the mean value theorem there are $\xi_n\in(1,\alpha_n)$ with the
bracket equal to $-B'_{\sigma_{\hat c,s}(\omega_n)}(\xi_n)
\ge-c_{\sigma_{\hat c,s}(\omega_n)}(\xi_n)$ by \eqref{eq:majorant}. Joint
continuity of $(\alpha,\sigma)\mapsto c_{\sigma}(\alpha)$
(Lemma~\ref{lem:b-derivative}) and
$\sigma_{\hat c,s}(\omega_n)\to\sigma_{\hat c,s}(\omega^{*})$ (the maps
$\cG_{\hat c,s}$ are continuous) give
$c_{\sigma(\omega_n)}(\xi_n)\to
c_{\sigma(\omega^{*})}(1)=B'_{\sigma(\omega^{*})}(1)$. Hence
\begin{equation*}
  \liminf_n\ q_n(\hat c)\,e_{\hat c}(\alpha_n,\omega_n)
  \;\ge\;\frac{q^{*}(\hat c)}{\ln2\,P_{\hat c}(\omega^{*})}
  \sum_{s}\big(-B'_{\sigma_{\hat c,s}(\omega^{*})}(1)\big)
  \;=\;\sum_{s}\Dtwo\big(\sigma_{\hat c,s}(\omega^{*})\,\big\|\,
  \cZ(\sigma_{\hat c,s}(\omega^{*}))\big),
\end{equation*}
using $q^{*}(\hat c)=P_{\hat c}(\omega^{*})$. Dropping the non-negative
divergence part and summing over $\hat c$ (a finite sum, so
$\liminf$ passes through),
$\liminf_n\hlow(\alpha_n)\ge h_1$-objective of
$\omega^{*}\ge h_1$. Combining with the upper bound proves
\eqref{eq:oracle}.
\end{proof}
